\section*{Data and code availability}

The benchmark definition and window index, the evaluation harness, the significance code, the figure scripts and
the aggregated score tables behind every number in this paper are available
at \url{https://github.com/ihpolash/BanglaWeatherBench}. The exact snapshot this paper describes is
archived as v1.0.2 at \url{https://doi.org/10.5281/zenodo.24358236}.
Source data are redistributed only where licensing permits: the BMD station record is CC-BY-4.0 \citep{zubair2024bmd}; GHCN-Daily and NASA POWER are public domain; CHIRPS is open \citep{funk2015chirps}. Derived tracks are rebuilt by the released pipeline rather than shipped, as are the per-window forecast cache and per-window scores, which run to several gigabytes. The release includes a datasheet in the
sense of \citet{gebru2021datasheets}, covering composition, collection, preprocessing, uses and limitations.

\section*{Reproducibility}

Model inference runs in per-model isolated environments; every forecast file is verified for row counts, window
identity, lead coverage, missing values, quantile monotonicity and non-negativity before it enters a leaderboard.
The paper can be reproduced at three depths, from cheapest to complete.
\begin{enumerate}
  \item \emph{Check the headline numbers} (seconds, CPU, fresh clone): \path{scripts/audit_claims.py} recomputes
    them from the committed score tables and window index, and \path{scripts/make_leaderboard_table.py}
    regenerates Tables~\ref{tab:leaderboard}, \ref{tab:evidence} and~\ref{tab:comparators} from the same files.
  \item \emph{Recompute the statistics} (CPU, after step 3's forecasts exist): \path{scripts/build_leaderboard.py}
    scores the per-window forecasts and \path{scripts/run_significance.py} produces every test and interval.
    \path{scripts/validation_selection.py} chooses the comparators on the validation forecasts and runs the primary
    family; \path{scripts/split_sample_check.py} and the test-chosen comparison come from the same score tables.
  \item \emph{Rerun everything}: \path{scripts/run_all.sh} rebuilds the data tracks and window index from the
    public sources. After that, the foundation models take 4.16 GPU-hours on a Tesla T4; the deep-learning
    baselines, one to two GPU-hours; the statistical and LightGBM baselines, about two hours on eight CPU cores; and
    AutoARIMA, about 20.7 CPU-hours. Setting \path{BWB_SPLIT=val} makes the same runners produce the
    2011--2015 validation forecasts that choose the comparators.
\end{enumerate}
\paragraph{Model versions.} Checkpoints and packages change over time, so a rerun may not match to the last
digit. We used the Hugging Face checkpoints \path{amazon/chronos-2}, \path{amazon/chronos-bolt-base},
\path{google/timesfm-2.5-200m-pytorch}, \path{Datadog/Toto-2.0-313m}, \path{NX-AI/TiRex},
\path{Salesforce/moirai-2.0-R-small}, \path{ibm-granite/granite-timeseries-ttm-r2} and
\path{thuml/sundial-base-128m}. The adapters load them without a pinned revision, so \path{envs/fm_kaggle/}
lists the commit of each repository that was current at run time; passing that commit as the revision
loads each checkpoint as it stood when we ran it.

\paragraph{Environments.} The local environments lock neuralforecast 3.2.2, statsforecast 2.0.1, LightGBM 4.7.0,
chronos-forecasting 2.3.2 and PyTorch 2.14.0. The GPU runs on the test split used the host's PyTorch 2.10.0 and
installed each foundation model's package separately: chronos-forecasting 2.3.2 with transformers 5.17.0
(Chronos-2, Chronos-Bolt), timesfm 3.0.2, toto-2 2.0.0, tirex-ts 1.4.2, uni2ts 2.0.0, granite-tsfm 0.3.9 with
transformers 5.17.0, and transformers 4.40.1 for Sundial. Those runs did not pin versions. The versions above come
from install logs of the same notebook on the day of the main run and of the later context runs, and
\path{envs/fm_kaggle/} turns them into one requirements file per foundation model.

The validation-period runs, made later on an updated Kaggle image, pin the same package versions. Five foundation
models ran there on Python 3.13 with PyTorch 2.11. Toto-2.0, Moirai-2 and Sundial do not install on that image
and ran on Python 3.12 with PyTorch 2.10, as on the test split. Every notebook is in \path{kaggle/}, and the
original logs of the Kaggle runs, test and validation, are in \path{reports/kaggle_logs/}.

\paragraph{Hardware and randomness.} The GPU runs used Kaggle's Tesla T4 GPUs, with PyTorch built for
CUDA 12.8; the five validation models on PyTorch 2.11 used its CUDA 13.0 build. Seven foundation models return
quantiles directly and involve no sampling. Sundial draws 100 sample paths with no fixed seed, so a rerun
reproduces its scores only up to sampling noise. The neural baselines, LightGBM and the assignment of gap patterns
to windows use seed 0, and the statistical baselines are deterministic. GPU training is still not bitwise
reproducible.

\paragraph{Data.} \path{reports/data_checksums.sha256} records SHA-256 hashes of every downloaded source file and
processed track as used, so a rebuild can check that it starts from the same data. The sources themselves, such as
NASA POWER, may be revised after our download.

\section*{Declaration on the use of generative AI}

The author used an AI assistant (Claude, Anthropic) for implementation support, analysis and figure scripting, and
for drafting and editing this manuscript. All of that work was carried out under the author's direction and
review. The research
questions, benchmark design, evaluation protocol and the interpretation of the results are the author's. All
numerical results were produced by the released code and were recomputed from the released artefacts by an audit
script before submission; every reference was verified by the author against its primary source. The author takes
full responsibility for all content of this paper, irrespective of how it was drafted.
